\documentclass[11pt]{article}
\usepackage[T1]{fontenc}
\usepackage{lmodern,amsmath,amssymb,amsthm}
\usepackage[margin=1in]{geometry}
\usepackage[hidelinks]{hyperref}
\newcommand{\Sha}{\mathop{\mathrm{Sha}}\nolimits}
\title{Splitting Genus-One Torsors and Rank Growth in Quadratic Towers}
\author{\small \href{mailto:admin@cosmosgeek.org}{admin@cosmosgeek.org}\qquad \href{mailto:b321014@iiit-bh.ac.in}{b321014@iiit-bh.ac.in}}
\date{Working draft for specialist review}
\begin{document}
\maketitle
\begin{abstract}
Let $E$ be an elliptic curve over a totally real number field $K$ with
irreducible two-torsion. A locally soluble torsor of period $2^k$ can
be split in a tower of $k$ totally real quadratic extensions, with
splitting prescribed at finitely many places. The differences between
conjugates of its splitting point generate the full rational augmentation
representation in the Mordell--Weil group of the normal closure. In
particular the rank increases at every step, and the gain in the normal
closure is at least $2^k-1$.
\end{abstract}
\medskip
\noindent\textbf{Theorem.}
Let $K$ be a totally real number field, and let
$E/K$ be an elliptic curve such that $E[2]$ is an
irreducible $\mathbb F_2[G_K]$-module. Suppose
\[
\alpha\in\Sha(E/K)
\]
has exact order $m=2^k$, with $k\ge1$.

For any prescribed finite set $S$ of finite places
of $K$, there exist a tower
\[
K=F_0\subset F_1\subset\cdots\subset F_k=F,
\qquad [F_i:F_{i-1}]=2,
\]
and a point $x$ on the $\alpha$-torsor over $F$,
with the following properties.

\begin{enumerate}
\item Every $F_i$ is totally real, and all places
above $S$ split completely at every step.

\item The restriction of $\alpha$ to $F_i$ has
exact order $2^{k-i}$.

\item The normal closure $L/K$ of $F/K$ is a
totally real Galois $2$-extension and
$E(L)[2]=0$.

\item If $G=\operatorname{Gal}(L/K)$ and
$H=\operatorname{Gal}(L/F)$, there is an injective
$\mathbb Q[G]$-homomorphism
\[
\mathbb Q[G/H]_{\mathrm{aug}}
\hookrightarrow E(L)\otimes\mathbb Q.
\]
Here the augmentation representation is the
kernel of the degree map on $\mathbb Q[G/H]$.

\item The ranks satisfy
\[
\operatorname{rank}E(F_i)
-\operatorname{rank}E(F_{i-1})\ge1
\quad(1\le i\le k)
\]
and
\[
\operatorname{rank}E(L)-\operatorname{rank}E(K)
\ge2^k-1.
\]
\end{enumerate}

\medskip
\noindent\textbf{Proof.}

\emph{Splitting an order-two multiple.}
Assume inductively that
\[
\alpha_i=\operatorname{res}_{F_i/K}(\alpha)
\]
has exact order $2^r$, where $r=k-i\ge1$.
Put
\[
\beta_i=2^{r-1}\alpha_i.
\]
Then $\beta_i$ has exact order two and belongs
to $\Sha(E/F_i)$.

Let $C_i/F_i$ represent $\beta_i$. By the
Cassels--Lichtenbaum period--index theorem
\cite[Theorem 2.7(ii)]{FisherVisibility}, its
index is two.

For completeness, the obstruction to representing
an $F_i$-rational geometric divisor class by an
actual $F_i$-rational line bundle lies in
$\operatorname{Br}(F_i)$. Since $C_i$ has a point
over every completion, this obstruction is locally
zero. The Brauer--Hasse--Noether injection makes
it globally zero. Thus $C_i$ possesses a line
bundle $\mathcal L_i$ of degree two.

Riemann--Roch gives $h^0(C_i,\mathcal L_i)=2$.
A degree-two line bundle on a genus-one curve is
base-point-free, so it defines a degree-two map
\[
f_i:C_i\longrightarrow\mathbb P^1_{F_i}.
\]

At any completion $F_{i,v}$, local solubility
gives a point of $C_i(F_{i,v})$. Since the branch
locus is finite and the set of local points is
infinite, we may choose a point away from it.

Its image under $f_i$ has a fibre consisting
of two distinct local points. By the real
implicit-function theorem or nonarchimedean
Hensel lifting, there is an open neighborhood
of that image whose fibres split into two
distinct $F_{i,v}$-points.

Apply weak approximation on $\mathbb P^1(F_i)$
to choose $t_i$ in these neighborhoods at every
real place and at every place above $S$, avoiding
the branch locus.

The fibre $f_i^{-1}(t_i)$ is a reduced scheme
of degree two. It cannot contain an
$F_i$-rational point, since $\beta_i\ne0$.
It is therefore the spectrum of a quadratic
field extension
\[
F_{i+1}/F_i.
\]

This extension is totally real and splits at
every specified finite place. It supplies a
point of $C_i$, so
\[
\operatorname{res}_{F_{i+1}/F_i}(\beta_i)=0.
\]
Consequently $\alpha_{i+1}$ has order at most
$2^{r-1}$.

On the other hand, restriction--corestriction
gives
\[
\operatorname{cor}_{F_{i+1}/F_i}(\alpha_{i+1})
=2\alpha_i,
\]
whose order is exactly $2^{r-1}$. Hence
\[
\operatorname{ord}(\alpha_{i+1})=2^{r-1}.
\]

Induction constructs the asserted tower.
At $i=k$, the class $\alpha_k$ is zero.
Therefore the original $\alpha$-torsor $C$
has an $F$-rational point $x$.

\medskip
\emph{The normal closure and its two-torsion.}
Let $L_i$ be the normal closure of $F_i/K$.
Inductively, $L_{i-1}/K$ has $2$-group Galois
group. Every conjugate of $F_i/F_{i-1}$ becomes
an extension of degree at most two after base
change to $L_{i-1}$.

Their compositum $L_i$ therefore has elementary
abelian $2$-group Galois group over $L_{i-1}$.
Thus $\operatorname{Gal}(L_i/K)$ is a $2$-group.
All conjugates of the totally real field $F_i$
are real, so $L_i$ is totally real.

In particular $L=L_k$ has the asserted
properties.

Since $E[2]$ is irreducible over $\mathbb F_2$,
the image of $G_K$ in
\[
\operatorname{GL}_2(\mathbb F_2)\simeq S_3
\]
is either $C_3$ or $S_3$.

The image of $G_L$ is normal and has index a
power of two in this image. It must therefore
contain the subgroup $C_3$.

A nonidentity element of $C_3$ fixes no
nonzero vector in $\mathbb F_2^2$. Hence
\[
E(L)[2]=0.
\]
In particular, every torsion point of $E(L)$
has odd order.

\medskip
\emph{The augmentation representation.}
The period of $C/K$ is $m$.
If $x$ were defined over an intermediate field
$F'/K$ of degree less than $m$, then
restriction--corestriction would give
\[
[F':K]\alpha=0,
\]
contradicting the exact order $m$.
Thus $K(x)=F$.

Write
\[
X=G/H,\qquad |X|=m.
\]
The conjugates $x_g=gx$, indexed by $X$, are
distinct and form one transitive $G$-set.

Let
\[
I_X=\ker\left(\mathbb Z[X]
\xrightarrow{\deg}\mathbb Z\right).
\]

The canonical identification
$\operatorname{Pic}^0(C_L)=E(L)$ gives a
$G$-equivariant homomorphism
\[
\Phi:I_X\longrightarrow E(L)/E(L)_{\rm tors},
\]
defined by
\[
\Phi\left(\sum_{g\in X}a_ge_g\right)
=
\left[\sum_{g\in X}a_g[x_g]\right]
\bmod E(L)_{\rm tors}.
\]
Every divisor in this formula has degree zero.

We prove that $\Phi$ is injective.

Suppose its kernel $J$ is nonzero.
The group $E(L)/E(L)_{\rm tors}$ is
torsion-free. Therefore $J$ is saturated in
$I_X$. Since $I_X$ itself is saturated in
$\mathbb Z[X]$, the composite inclusion gives
\[
J/2J\hookrightarrow
\mathbb F_2[X].
\]

The left side is a nonzero
$\mathbb F_2[G]$-module. Every nonzero module
over a finite $2$-group in characteristic two
has nonzero invariants: the augmentation ideal
of $\mathbb F_2[G]$ is nilpotent, and its
highest nonzero action on the module produces
a nonzero vector annihilated by that ideal.

Hence $(J/2J)^G\ne0$.

Because $G$ acts transitively on $X$,
\[
\mathbb F_2[X]^G
=
\mathbb F_2\left(\sum_{g\in X}e_g\right).
\]

It follows that $J$ contains an element
\[
z=\sum_{g\in X}(1+2a_g)e_g.
\]
Since $z\in I_X$,
\[
m+2\sum_g a_g=0,
\qquad
\sum_g a_g=-m/2.
\]

Define the divisors
\[
D=\sum_g[x_g],
\qquad
A=-\sum_g a_g[x_g].
\]
Their degrees are $m$ and $m/2$, respectively.

The divisor $D$ is defined over $K$.
Because $z\in\ker\Phi$, the degree-zero
class
\[
T=[D-2A]\in E(L)
\]
is torsion.

Every torsion point of $E(L)$ has odd order.
Multiplication by two is consequently an
automorphism of the finite torsion subgroup.
Choose $S\in E(L)_{\rm tors}$ with $2S=T$.

Now put
\[
B=[A]+S\in\operatorname{Pic}^{m/2}(C_L)(L).
\]
Then
\[
2B=[D].
\]

For every $g\in G$,
\[
2(gB-B)=g[D]-[D]=0.
\]
Since $E(L)[2]=0$, we have $gB=B$.

Thus $B$ descends to a $K$-rational point
of the degree-$m/2$ Picard \emph{scheme}:
\[
B\in\operatorname{Pic}^{m/2}_{C/K}(K).
\]

This distinction is essential. We do not
claim that $B$ necessarily lifts to an actual
$K$-rational divisor. The period is defined
using $K$-rational geometric divisor classes,
so a rational point of this Picard scheme is
exactly what is needed.

The existence of $B$ contradicts the
period $m$ of $C$.

Therefore $\ker\Phi=0$, and tensoring with
$\mathbb Q$ proves
\[
\mathbb Q[X]_{\rm aug}
\hookrightarrow
E(L)\otimes\mathbb Q.
\]

\medskip
\emph{Rank at each step.}
Put
\[
H_i=\operatorname{Gal}(L/F_i),
\qquad
o_i=\#(H_i\backslash X).
\]

For a permutation representation,
\[
\dim_{\mathbb Q}
\left(\mathbb Q[X]_{\rm aug}\right)^{H_i}
=o_i-1.
\]

Because $H_i$ is a normal subgroup of index
two in $H_{i-1}$, each $H_{i-1}$-orbit on $X$
either remains one orbit or splits into two
$H_i$-orbits.

The distinguished orbit containing $H\in X$
necessarily splits: it is identified with
$H_{i-1}/H$, and $H\subseteq H_i$.

Consequently
\[
o_i-o_{i-1}\ge1.
\]

By Maschke's theorem, the injection proved
above splits as a map of rational
$G$-representations. Its complementary
representation can only gain invariant
vectors as $H_{i-1}$ is replaced by $H_i$.

Using Galois descent,
\[
\begin{aligned}
\operatorname{rank}E(F_i)
-\operatorname{rank}E(F_{i-1})
&\ge o_i-o_{i-1}\\
&\ge1.
\end{aligned}
\]

Finally, $G$ has one orbit on $X$, whereas
the trivial subgroup has $m$ orbits. Therefore
\[
\operatorname{rank}E(L)
-\operatorname{rank}E(K)\ge m-1.
\]
This completes the proof.
\hfill$\square$

\medskip
\noindent\textbf{Scope and literature comparison.}
Cassels--Lichtenbaum supplies the period--index
input \cite{FisherVisibility}. Radi\v{c}evi\'c
studies degree-period capitulation with
discriminant bounds in a different setting
\cite{Radicevic}. Rank growth over nonabelian
extensions is studied by Mazur and Rubin
\cite{MazurRubin}. These references do not,
by themselves, establish that the full
augmentation statement above is new.

The rank gain concerns $E$ after extending its ground field. The original
central order of $L(E/K,s)$ is a separate question.


\appendix
\section{Heights of representatives over splitting fields}
\label{sec:heights}

For an elliptic curve $E$ over a number field $k$ write
$\|X\|^2=\widehat h(X)$ for the absolute N\'eron--Tate height,
extended as a quadratic form on $E(F)\otimes_{\mathbb Z}\mathbb R$.
For a positive integer $q$ let $\delta_{q,F}:E(F)/qE(F)\to
H^1(F,E[q])$ be the Kummer map. The following estimate concerns the
part of a representative perpendicular to $E(k)\otimes\mathbb R$.

\medskip
\noindent\textbf{Proposition A.1.}
Let $q=p^m$ for a prime $p$, let $F/k$ be a finite Galois extension
with $E(F)[p]=0$, and set $G=\operatorname{Gal}(F/k)$ and
$e=|G|^{-1}\sum_{g\in G}g$. Suppose that $\kappa\in H^1(k,E[q])$
restricts to $\delta_{q,F}(P)$ and that its image $\xi\in H^1(k,E)$
has exact order $p^a>1$. Let $\mu>0$ be a lower bound for the height
of every nontorsion point in $E(F)$. Then for every $Q\in E(F)$,
\begin{equation}\label{eq:transverse-bound}
 \frac{\|(1-e)(P-qQ)\|^2}{q^2}
 \ \ge\ \frac{\mu}{2}(1-p^{-a}).
\end{equation}

\begin{proof}
Put $P'=P-qQ$ and let $d$ be the size of the $G$-orbit of $P'$
modulo $T=E(F)_{\rm tors}$. Take $K$ to be the fixed field of the
stabilizer, so $[K:k]=d$. The integer $D=|T|$ is prime to $q$,
and $DP'\in E(K)$. Restriction from $K$ to $F$ is injective on
$H^1(K,E[q])$ by inflation--restriction and $E(F)[q]=0$. The classes
$D\,\operatorname{res}_{K/k}\kappa$ and $\delta_{q,K}(DP')$ agree
over $F$, hence over $K$. It follows that $D\,\operatorname{res}_{K/k}
\xi=0$. Since $\xi$ is $p$-primary, $K$ splits $\xi$.
Corestriction now gives $p^a\mid d$.

The class of $P'$ in $E(F)/qE(F)$ is $G$-invariant by injectivity
and equivariance of the Kummer map. Each $gP'-P'$ therefore equals
$q$ times a point of $E(F)$. A fraction $1-1/d$ of these differences
are nontorsion and have height at least $q^2\mu$. The action preserves
the height pairing, and $e$ is orthogonal projection to the fixed
space. Consequently
\[
 \frac1{|G|}\sum_{g\in G}\|gP'-P'\|^2
 =2\|P'\|^2-2\langle eP',P'\rangle
 =2\|(1-e)P'\|^2 .
\]
The two inequalities $d\ge p^a$ and
$2\|(1-e)P'\|^2\ge(1-1/d)q^2\mu$ prove the proposition.
\end{proof}

\noindent\textbf{Proposition A.2.}
Fix $\kappa$ and its nonzero image $\xi$ as above. Let
$F_i/k$ be pairwise linearly disjoint finite Galois extensions,
with $E(F_i)[p]=0$ and $\operatorname{res}_{F_i/k}\kappa
=\delta_{q,F_i}(P_i)$. Assume a common lower bound $\mu>0$ for
the heights of nontorsion points in all the $F_i$ and $F_iF_j$.
Define
\[
 e_i=\frac1{|\operatorname{Gal}(F_i/k)|}
        \sum_{g\in\operatorname{Gal}(F_i/k)}g,\qquad
 u_i=(1-e_i)P_i,\qquad \beta_i=\|u_i\|^2/q^2.
\]
Put $\Lambda=E(k)/E(k)_{\rm tors}$, $V_0=\Lambda\otimes\mathbb R$,
and let $x_i$ be the class of $e_iP_i/q$ in $V_0/\Lambda$. Then
\begin{equation}\label{eq:torus-distance}
 \operatorname{dist}(x_i,x_j)^2+\beta_i+\beta_j\ \ge\ \mu
 \qquad (i\ne j).
\end{equation}
In particular, if $0\le\eta<\mu/2$, the number of $i$
for which $\beta_i\le\eta$ is at most
\[
 \frac{2^r\sqrt{\operatorname{Reg}(E/k)}}
 {v_r(\mu-2\eta)^{r/2}},
 \qquad r=\operatorname{rank}E(k),
\]
where the regulator uses the given norm and $v_r$ is the Euclidean
unit-ball volume. For $r=0$ this bound is one. Along any infinite
sequence of such fields, $\liminf_i\beta_i\ge\mu/2$.

\begin{proof}
Over $F_iF_j$ both $P_i$ and $P_j$ represent the same Kummer class.
Write $P_i-P_j=qR_{ij}$ for $R_{ij}\in E(F_iF_j)$. The vector
$u_i$ is nonzero by Proposition~A.1. Averaging the pairing with $u_j$
over $\operatorname{Gal}(F_iF_j/F_j)$ shows $u_i\perp u_j$,
since this group restricts onto $\operatorname{Gal}(F_i/k)$ and
fixes $u_j$. Both are orthogonal to $V_0$. For $Q\in E(k)$ the point
$R_{ij}-Q$ is nontorsion since its component $(u_i-u_j)/q$ is nonzero.
Therefore
\[
 \mu\le \|R_{ij}-Q\|^2
 =\|(e_iP_i-e_jP_j)/q-Q\|^2+\beta_i+\beta_j.
\]
Minimization over $Q$ gives \eqref{eq:torus-distance}.

If both $\beta_i,\beta_j\le\eta$, their torus distance is at least
$\rho=\sqrt{\mu-2\eta}$. In positive rank each nonzero lattice
vector has norm at least $\sqrt{\mu}\ge\rho$. Hence the
\emph{open} balls of radius $\rho/2$ project injectively to the quotient and retain their Euclidean volume.
Indeed, two points in one open
ball have distance strictly less than $\rho$ in the covering space,
so their difference cannot be a nonzero period. Distinct balls in
the quotient have disjoint interiors. Their volumes are therefore
at most the total covolume $\sqrt{\operatorname{Reg}(E/k)}$, giving
the claimed bound. In rank zero the torus consists of one point,
and \eqref{eq:torus-distance} directly excludes a second index with
$\beta_i\le\eta$. Applying the bound for each $\eta<\mu/2$ gives
the limit inferior.
\end{proof}

For a single number field $F$ the positive height minimum follows
from the Mordell--Weil lattice. In an infinite family the common
bound is a substantive hypothesis; Silverman supplies it when all
fields lie in $k^{\rm ab}$ \cite{SilvermanHeight}. The argument
allows all classes in the Weil--Ch\^atelet group, including locally
soluble torsors. An upper height estimate for particular representatives
would require separate arithmetic information.

\begin{thebibliography}{9}

\bibitem{FisherVisibility}
T.~Fisher,
\emph{Visibility of Tate--Shafarevich Groups},
Warwick lecture notes (2008),
Theorem~2.7(ii).
\url{https://www.dpmms.cam.ac.uk/~taf1000/papers/visibility.pdf}

\bibitem{Radicevic}
L.~Radi\v{c}evi\'c,
\emph{Capitulation Discriminants of Genus One Curves},
arXiv:2212.00838 (2022).
\url{https://arxiv.org/abs/2212.00838}

\bibitem{MazurRubin}
B.~Mazur and K.~Rubin,
\emph{Growth of Selmer Rank in Nonabelian Extensions
of Number Fields},
arXiv:math/0703363.
\url{https://arxiv.org/abs/math/0703363}

\bibitem{SilvermanHeight}
J. H. Silverman, \emph{A lower bound for the canonical height on
elliptic curves over abelian extensions}, J. Number Theory 104
(2004), 353--372. \url{https://arxiv.org/abs/math/0305041}

\end{thebibliography}

\end{document}
